\bookchapter{Chapter 51: Building the UI — Getting Started with Jetpack Compose}

\chapterintro{The UI is everything the user sees. Before Compose came along, Android UIs were written in XML — as tedious as fill-in-the-blank exercises. In 2021 Google introduced Jetpack Compose: you “draw” the UI directly in Kotlin code, what you see is what you get, and the preview refreshes the moment you change a line of code. By 2026, Compose is the clear mainstream for new projects. This chapter is where we learn to “draw.”}

\booksection{I. Compose's Big Idea: The UI Is a Function}

Compose's core idea in one sentence: UI = function. Everything that can appear on screen (text, buttons, images) is a “composable function,” marked with @Composable. Nest them like building blocks, and that's your UI.

\begin{lstlisting}[language=Kotlin]
@Composable
fun Greeting(name: String) {
    Text(text = "Hello, $name!")     // a "text building block"
}

// Nested composition: big blocks around small ones
@Composable
fun HomeScreen() {
    Column {                          // vertical stack
        Greeting(name = "Alex")
        Greeting(name = "Sam")
        Button(onClick = { /* what happens on tap */ }) {
            Text("Tap me")
        }
    }
}
\end{lstlisting}

Note the @Composable “hat” — it tells the compiler “this is a UI building block.” Blocks can nest infinitely; that's the whole secret of Compose.

\booksection{II. Basic Components: Text, Button, Input Field, Image}

Compose's components are “building blocks.” The four most common:

\begin{center}
\begin{tabularx}{\linewidth}{|X|X|X|}
\hline
\textbf{Component} & \textbf{What It Does} & \textbf{Required Parameter} \tabularnewline
\hline
Text & Displays text & text (the text to show) \tabularnewline
Button & Button & onClick (the click callback) \tabularnewline
TextField & TextField & value and onValueChange (the callback for input changes) \tabularnewline
Image & Displays images & painter (the image resource) \tabularnewline
\hline
\end{tabularx}
\end{center}

\begin{center}\small\color{gray}Compose's Four Basic Components\end{center}

\begin{lstlisting}[language=Kotlin]
@Composable
fun BasicWidgets() {
    Column {
        Text(text = "This is some text")
        Button(onClick = { println("clicked") }) {
            Text("Button")
        }
        var input by remember { mutableStateOf("") }
        TextField(
            value = input,
            onValueChange = { input = it },
            label = { Text("Type something") }
        )
    }
}
\end{lstlisting}

In those two lines for the input field you'll spot remember \{ mutableStateOf(“”) \} — that's Compose's “state” magic, covered in detail in Section 5 of this chapter. For now, just get acquainted with it.

\booksection{III. The Layout Trio: Column, Row, Box}

You have the components — how do you arrange them? Remember three “boxes”:

\begin{itemize}
  \item Column: stacks things vertically (like a line at the cafeteria counter);
  \item Row: lines things up horizontally (like sitting side by side);
  \item Box: stacks things on top of each other (like a pile of playing cards — great for badges and overlays).
\end{itemize}

\begin{lstlisting}[language=Kotlin]
Column {
    Text("First row")
    Row {
        Text("Left")
        Text("Right")
    }
    Box {
        Text("Bottom layer")
        Text("Corner badge stacked on top")
    }
}
\end{lstlisting}

Want the layout to look better? Add a Modifier to components — it's like putting makeup on your building blocks:

\begin{lstlisting}[language=Kotlin]
Text(
    text = "Pretty text",
    modifier = Modifier
        .padding(16.dp)          // padding on all sides
        .fillMaxWidth()          // fill the width
        .background(Color.Blue)  // background color
        .clickable { /* tap */ }// clickable
)
// Chained calls: one after another, order matters (padding first, then background — the background excludes the padding)
\end{lstlisting}

\begin{tcbwarn}{Note: Modifier Order Matters}
padding(16.dp).background(Blue) and background(Blue).padding(16.dp) give different results — the first is “spacing first, then paint” (the background excludes the padding), the second is “paint first, then spacing” (the background includes the padding). This is where beginners most often paint themselves a “mysterious white border.” When unsure, swap the order — the preview instantly tells you the answer.

\end{tcbwarn}

\booksection{IV. State: Bringing the UI to Life}

A static UI is boring. To make a number go up by 1 when a button is clicked, or let an input field accept typing, you need state. Compose's rule: the UI always faithfully reflects the state — when the state changes, the UI redraws automatically.

\begin{lstlisting}[language=Kotlin]
@Composable
fun Counter() {
    // remember + mutableStateOf: remember this changing value
    var count by remember { mutableStateOf(0) }

    Column {
        Text(text = "You tapped $count times")
        Button(onClick = { count++ }) {   // change the state
            Text("Tap me +1")                // the UI refreshes automatically
        }
    }
}
\end{lstlisting}

Three lines of code, one live counter. This is where Compose feels better than the old way: no manually “refreshing the UI” — just change the data. Remember the two key functions: remember (keeps a value from being wiped when the UI recomposes) and mutableStateOf (tells Compose “redraw when this value changes”).

User taps

State changes (count++)

Compose notices the change

UI redraws automatically (recomposition)

\begin{center}\small\color{gray}Figure 3: Compose's one-way data flow from state to UI\end{center}

\booksection{V. Theming and Dark Mode: Configure Once, Apply Everywhere}

The MyFirstAppTheme \{ ... \} in the template project is the theme. Based on the Material 3 design spec, it manages colors, fonts, and shapes in one place. Supporting dark mode takes only three steps:

\begin{enumerate}
  \item Define both a light and a dark color scheme in the theme (the template already generates the ColorScheme);
  \item When the system switches to dark mode, isSystemInDarkTheme() detects it automatically;
  \item Use the colors the theme provides (like MaterialTheme.colorScheme.primary) instead of hard-coded color values.
\end{enumerate}

\begin{tcbinfo}{Tip: Material 3 Dynamic Color}
Starting with Android 12, “Material You” lets an app pull colors from your wallpaper. Compose templates support it by default (dynamicColor). Your app will “change outfits” on different users' phones — that's the romance of modern Android.

\end{tcbinfo}

\booksection{VI. Preview, Animation, and Live Edit: Compose's “Cheat Codes”}

Three “cheat codes” for Compose's developer experience — once you use them, there's no going back:

\begin{itemize}
  \item @Preview: add the @Preview annotation to a Composable and Android Studio instantly renders the UI on the right side — what you see is what you get as you edit, no need to install on a phone every time;
  \item Live Edit: the 2026 Android Studio supports "hot updating" certain changes on the emulator or a real device — small tweaks like changing text color appear on the device in seconds;
  \item Animation APIs: animateFloatAsState and AnimatedVisibility let elements transition smoothly in a few lines of code — that "premium feel" is often just these few dozen lines away.
\end{itemize}

\begin{lstlisting}[language=Kotlin]
@Composable
fun AnimatedBox() {
    var big by remember { mutableStateOf(false) }
    // when the value changes, it smoothly animates to the target
    val size by animateDpAsState(if (big) 120.dp else 60.dp)

    Box(Modifier
        .size(size)
        .background(MaterialTheme.colorScheme.primary)
        .clickable { big = !big }
    )
}
\end{lstlisting}

\begin{tcbinfo}{Tip: Function First, Animation Later}
A common beginner mistake is fussing over animations before the UI even works. Get the feature running first, then add motion at the end — that's the Simplicity First principle too: animations are the icing on the cake, features are the cake itself.

\end{tcbinfo}

\begin{tcbpractice}{Hands-On: Build a "Mood Button"}
Turn the template project into a mood switcher: one button that, when tapped, cycles through happy sad angry, with the text and background color changing along. Hint: use when to pick the wording and color based on the state. Send it to friends when it's done and collect your first round of “whoa, you can make apps now”.

\end{tcbpractice}

\begin{tcbsummary}{Chapter Summary}
\begin{itemize}
  \item Compose's idea: the UI = @Composable functions, nested like building blocks
  \item Four basic components: Text / Button / TextField / Image
  \item The layout trio: Column stacks vertically, Row lines up horizontally, Box stacks on top
  \item Modifier chains apply "makeup" in sequence: padding, fillMaxWidth, background, clickable — order matters
  \item State is the soul: remember + mutableStateOf — change the data and the UI redraws itself
  \item Theming with Material 3 + dynamic color: configure dark mode once and it applies everywhere
\end{itemize}

\end{tcbsummary}
